% ===== PRÉAMBULE CANONIQUE — à coller VERBATIM en tête de CHAQUE .tex =====
\documentclass[11pt,a4paper]{article}
\usepackage[utf8]{inputenc}\usepackage[T1]{fontenc}\usepackage{lmodern}
\usepackage[french]{babel}
\usepackage{amsmath,amssymb,mathtools}\usepackage{icomma}
\usepackage[version=4]{mhchem}          % \ce{...} : équations & formules
\usepackage{siunitx}\sisetup{locale=FR} % \qty{}{}, \si{}, \ang{}
\usepackage{chemfig}                    % \chemfig{...} : structures développées
\usepackage{xcolor}\usepackage{colortbl}
\usepackage{tikz}\usepackage{pgfplots}\pgfplotsset{compat=1.18}
\usepackage[most]{tcolorbox}\usepackage{enumitem}\usepackage{array,booktabs}
\usepackage[a4paper,margin=2.2cm]{geometry}
\usetikzlibrary{arrows.meta,calc,angles,quotes,positioning,patterns,backgrounds,shapes.geometric}
\definecolor{primary}{RGB}{222,49,99}\definecolor{primarydark}{RGB}{150,22,55}
\definecolor{defcol}{RGB}{0,114,178}\definecolor{methcol}{RGB}{52,92,148}
\definecolor{excol}{RGB}{0,135,90}\definecolor{attncol}{RGB}{200,40,55}
\tcbset{boxbase/.style={enhanced,breakable,boxrule=0.9pt,arc=2.5pt,
  left=3mm,right=3mm,top=2.3mm,bottom=2.3mm,fonttitle=\bfseries,coltitle=white}}
% --- boîtes TUTORIEL (noms EXACTS, ne pas renommer) ---
\newtcolorbox{cle}[1][]{boxbase,colback=defcol!6,colframe=defcol,
  title={\textsc{À retenir}\ifx\relax#1\relax\relax\else\textmd{ : \itshape #1}\fi}}
\newtcolorbox{methode}[1][]{boxbase,colback=methcol!7,colframe=methcol,sharp corners,
  title={\textsc{Méthode}\ifx\relax#1\relax\relax\else\textmd{ --- \itshape #1}\fi}}
\newtcolorbox{exemplebox}[1][]{boxbase,colback=excol!5,colframe=excol,
  title={\textsc{Exemple}\ifx\relax#1\relax\relax\else\textmd{ : \itshape #1}\fi}}
\newtcolorbox{piege}[1][]{boxbase,colback=attncol!6,colframe=attncol,
  title={\textsc{Piège}\ifx\relax#1\relax\relax\else\textmd{ : \itshape #1}\fi}}
\newtcolorbox{remarque}[1][]{boxbase,colback=black!4,colframe=black!45,
  title={\textsc{Remarque}\ifx\relax#1\relax\relax\else\textmd{ : \itshape #1}\fi}}
% --- boîtes EXERCICE / CORRIGÉ (noms EXACTS, auto-comptées) ---
\newtcolorbox[auto counter]{exo}[1][]{boxbase,colback=primary!4,colframe=primary,
  title={\textsc{Exercice}~\thetcbcounter\ifx\relax#1\relax\relax\else\textmd{ --- \itshape #1}\fi}}
\newtcolorbox[auto counter]{corrige}[1][]{boxbase,colback=excol!4,colframe=excol,
  title={\textsc{Corrigé}~\thetcbcounter\ifx\relax#1\relax\relax\else\textmd{ --- \itshape #1}\fi}}
\newcommand{\R}{\mathbb{R}}\newcommand{\N}{\mathbb{N}}\newcommand{\Z}{\mathbb{Z}}\newcommand{\ds}{\displaystyle}
% (un \usepackage{fancyhdr}+en-tête est possible ; il est ignoré côté web)
% =========================================================================
\begin{document}

\begin{corrige}[deux ions de valence 1]
Rapport $1:1$, aucun indice à ajouter.
\begin{enumerate}[label=\alph*)]
  \item \ce{NaCl}.
  \item \ce{KNO3}.
  \item \ce{NH4Cl}.
  \item \ce{AgBr}.
\end{enumerate}
\end{corrige}

\begin{corrige}[croiser une valence 1 et une valence 2]
La valence de l'un devient l'indice de l'autre.
\begin{enumerate}[label=\alph*)]
  \item \ce{CaCl2} (deux \ce{Cl-} pour un \ce{Ca^2+}).
  \item \ce{Na2O} (deux \ce{Na+} pour un \ce{O^2-}).
  \item \ce{K2S}.
  \item \ce{AlCl3} (trois \ce{Cl-} pour un \ce{Al^3+}).
\end{enumerate}
\end{corrige}

\begin{corrige}[penser à simplifier]
Après croisement, on divise les indices par leur diviseur commun.
\begin{enumerate}[label=\alph*)]
  \item \ce{Mg2O2} $\to$ \ce{MgO} (diviser par 2).
  \item \ce{Ca2S2} $\to$ \ce{CaS}.
  \item \ce{Al3(PO4)3} $\to$ \ce{AlPO4} (diviser par 3).
\end{enumerate}
\end{corrige}

\begin{corrige}[mettre l'ion polyatomique entre parenthèses]
On croise, puis on entoure l'anion polyatomique de parenthèses s'il est présent $\geq 2$ fois.
\begin{enumerate}[label=\alph*)]
  \item \ce{Mg(OH)2}.
  \item \ce{Ca(NO3)2}.
  \item \ce{(NH4)2SO4} (le cation polyatomique \ce{NH4+} est aussi parenthésé).
\end{enumerate}
\end{corrige}

\begin{corrige}[quand une valence 3 entre en jeu]
\begin{enumerate}[label=\alph*)]
  \item \ce{Al2O3} (croiser 3 et 2 : $\ce{Al2O3}$, non simplifiable).
  \item \ce{FeCl3} (pas de parenthèses : \ce{Cl} est monoatomique).
  \item \ce{Ca3(PO4)2} (croiser 2 et 3 : trois \ce{Ca^2+}, deux \ce{PO4^3-}).
\end{enumerate}
\end{corrige}

\end{document}
